\documentclass[conference]{IEEEtran}
\usepackage{cite}
\usepackage{amsmath,amssymb,amsfonts}
\usepackage{algorithmic}
\usepackage{graphicx}
\usepackage{textcomp}
\usepackage{xcolor}
\usepackage{booktabs}
\usepackage{array}
\usepackage{url}

\graphicspath{{figures/}{../../write_up/figures/}}

\begin{document}

\title{Edge-Feasible Intention Recognition and Spatial Navigation for Autonomous Mobile Robots in Shared Workspaces}

\author{
\IEEEauthorblockN{First Author\IEEEauthorrefmark{1}, Co-Author\IEEEauthorrefmark{1}, and Academic Supervisor\IEEEauthorrefmark{1}}
\IEEEauthorblockA{\IEEEauthorrefmark{1}Department of Computer Engineering, Faculty of Engineering\\
Ghana Communication Technology University (GCTU), Tesano, Accra, Ghana\\
Email: \{student, supervisor\}@gctu.edu.gh}
}

\maketitle

\begin{abstract}
Collaborative autonomous mobile robots (AMRs) operating in human-populated industrial and service environments require responsive, intuitive communication channels while maintaining deterministic safety. Existing vision-based intention recognition frameworks predominantly depend on power-intensive desktop GPUs or remote cloud infrastructure, introducing unpredictable network latency, vulnerability to wireless dropouts, and spatial-blind execution. This paper presents an edge-deployed, real-time human intention recognition and navigation architecture implemented natively on an embedded single-board computer (Raspberry Pi 5) without external GPU acceleration. The perception pipeline couples monocular person tracking with a scale- and distance-invariant 19-dimensional geometric hand feature extractor, feeding an edge-optimized Multilayer Perceptron (MLP) gesture classifier alongside an LSTM kinematic motion predictor. A centralized ROS~2 decision node (\texttt{brain\_node}) integrates rolling temporal majority voting, active subject-locking, and visual servoing with 2D LiDAR Simultaneous Localization and Mapping (SLAM). Physical trials across 180 live human-robot interaction runs demonstrated a 96.67\% aggregate classification accuracy across six operational commands under varying distances (1.00\,m to 2.50\,m) and ambient lighting conditions. The end-to-end perception-to-actuation cycle measured 132.0\,ms, strictly satisfying the 150\,ms human-interaction reaction threshold and complying with ISO 15066 collaborative robotic safety standards.
\end{abstract}

\begin{IEEEkeywords}
Human-Robot Interaction (HRI), Autonomous Mobile Robots, Edge Computing, ROS 2, Gesture Recognition, SLAM, Collaborative Safety.
\end{IEEEkeywords}

\section{Introduction}
Robotics engineering has transitioned rapidly from isolated industrial cages toward collaborative spaces where autonomous mobile robots share physical environments with human personnel. In these dynamic settings, intuitive Human-Robot Interaction (HRI) is essential for operational efficiency and worker safety. Non-verbal communication modalities—specifically natural hand gestures and body movement cues—provide an expressive, touchless interface suitable for acoustically noisy manufacturing floors or sterile healthcare environments where physical contact must be strictly avoided.

Despite recent advancements in computer vision, deploying vision-based intention recognition onto resource-constrained mobile robots presents three fundamental engineering bottlenecks:
\begin{enumerate}
    \item \textbf{GPU and Cloud Dependency:} Conventional deep neural networks often rely on power-hungry workstation GPUs or offboard cloud services \cite{mahmud2022}. In mobile robotics, reliance on remote servers introduces unpredictable network latency, high power consumption, and single-point communication failure risks that violate real-time safety thresholds \cite{tsitos2022}.
    \item \textbf{Spatial Coordinate Domain Shift:} Many vision models classify raw pixel coordinates of skeletal landmarks. Because coordinate values scale inversely with distance, unnormalized models suffer severe accuracy degradation as human operators traverse natural working distances (1.0\,m to 2.5\,m).
    \item \textbf{Spatial-Blind Actuation:} Existing gesture-controlled mobile prototypes frequently translate classified gestures directly into motor velocity commands without active environmental mapping or obstacle verification, creating collision hazards in cluttered indoor facilities \cite{li2023}.
\end{enumerate}

To overcome these challenges, this paper presents a unified, edge-deployed intention recognition and navigation architecture operating entirely on an embedded edge computing platform. The primary contributions of this work are:
\begin{itemize}
    \item A mathematically scale-invariant 19-dimensional hand feature representation that decouples landmark geometry from camera distance and operator hand morphology.
    \item A balanced 6,000-sample onboard gesture dataset and lightweight Multilayer Perceptron (MLP) inference pipeline achieving 99.1\% offline and 96.67\% physical real-time accuracy on edge CPU hardware.
    \item A multi-node ROS~2 architecture integrating active subject-locking, temporal voting, visual servoing, and 2D LiDAR SLAM, validated with an end-to-end latency budget of 132.0\,ms on physical robot hardware.
\end{itemize}

\section{Related Work}
Human intention inference in collaborative robotics spans kinematic trajectory forecasting, 3D gesture tracking, and intention-aware navigation.

Tsitos et al. \cite{tsitos2022} evaluated upper-limb kinematic intention prediction during a competitive human-robot reaching game using a single RGB-D sensor and OpenPose. By isolating wrist joint velocity profiles, their system predicted intention within strict latency bounds to govern an industrial UR3 manipulator. However, their architecture was restricted to a static tabletop manipulator without mobile base locomotion or explicit gesture command channels.

Mahmud et al. \cite{mahmud2022} developed a 3D skeletal tracking framework using a Microsoft Kinect depth sensor to classify pointing vectors and dynamic gestures across 24 human subjects. While achieving high classification accuracy, the pipeline required specialized, high-power depth hardware and validated navigation strictly within virtual simulation environments without addressing real-world wheel slip or embedded CPU constraints.

Li et al. \cite{li2023} formulated an intention-aware motion planning framework that modified ROS~2 Nav2 costmaps based on predicted human worker clearance behaviors. Although demonstrating reduced path replanning, their system was entirely passive (workers could not issue explicit directives) and remained unverified on physical edge hardware.

Recent lightweight ROS~2 vision architectures (e.g., Muhtadin et al. \cite{muhtadin2025}) demonstrated MediaPipe landmark classification, but evaluated performance solely on stationary industrial manipulators without mobile spatial mapping. In contrast, this paper resolves this technological divide by integrating explicit gesture commanding, kinematic motion prediction, and LiDAR SLAM navigation natively onto an embedded mobile testbed.

\section{System Architecture and Methodology}

\subsection{Physical Testbed and Hardware Design}
The experimental platform is a differential-drive autonomous mobile robot engineered for indoor navigation. The hardware consists of:
\begin{itemize}
    \item \textbf{Primary Edge Compute:} Raspberry Pi 5 single-board computer featuring a quad-core ARM Cortex-A76 processor (2.4\,GHz) and 8\,GB LPDDR4X RAM running Ubuntu 24.04 LTS and ROS~2 Humble.
    \item \textbf{Visual Perception:} Monocular 2MP USB camera ($640 \times 480$ resolution at 20\,FPS) mounted on a 2-DOF active pan-tilt servo gimbal.
    \item \textbf{Planar Range Sensing:} MS200 Time-of-Flight 2D LiDAR delivering $360^\circ$ planar range scans at 12.6\,Hz.
    \item \textbf{Low-Level Actuation Bridge:} Yahboom MicroROS expansion board powered by an Espressif ESP32-S3 microcontroller, communicating with the edge host via a 921,600 baud UART serial bus.
\end{itemize}

\subsection{Scale-Invariant Geometric Feature Formulation}
To eliminate spatial distortion across varying human-to-robot distances, raw 21-point MediaPipe hand landmarks $(x_k, y_k, z_k)$ are mapped into a 19-dimensional scale- and position-invariant geometric feature vector $\mathbf{f} \in \mathbb{R}^{19}$:

\textbf{1) Digit Curl Angles ($\mathbf{f}_{1\dots5}$):} For each of the five digits $i \in \{1,\dots,5\}$, the articulation curl angle $\theta_i$ is computed from the vector dot product between the proximal vector $\mathbf{u}_i = \mathbf{p}_{\text{mid}} - \mathbf{p}_{\text{base}}$ and distal vector $\mathbf{v}_i = \mathbf{p}_{\text{tip}} - \mathbf{p}_{\text{mid}}$:
\begin{equation}
\theta_i = \arccos\left( \frac{\mathbf{u}_i \cdot \mathbf{v}_i}{\|\mathbf{u}_i\| \|\mathbf{v}_i\| + \epsilon} \right), \quad i \in \{1,\dots,5\}
\end{equation}
where $\epsilon = 10^{-7}$ prevents division by zero. Because $\theta_i$ depends exclusively on angular vector orientation, it is mathematically invariant to image scale, camera distance, and uniform coordinate translation.

\textbf{2) Palm-Normalized Fingertip Distances ($\mathbf{f}_{6\dots10}$):} Euclidean distances from the wrist $\mathbf{p}_0$ to each fingertip $\mathbf{p}_{\text{tip}, i}$ are normalized by the instantaneous palm reference width $D_{\text{palm}} = \|\mathbf{p}_5 - \mathbf{p}_{17}\|$:
\begin{equation}
d_{\text{wrist}, i} = \frac{\|\mathbf{p}_{\text{tip}, i} - \mathbf{p}_0\|}{D_{\text{palm}} + \epsilon}, \quad i \in \{1,\dots,5\}
\end{equation}

\textbf{3) Inter-Fingertip Distances ($\mathbf{f}_{11\dots14}$):} Adjacent fingertip separations are normalized by palm width:
\begin{equation}
d_{\text{adj}, j} = \frac{\|\mathbf{p}_{\text{tip}, j+1} - \mathbf{p}_{\text{tip}, j}\|}{D_{\text{palm}} + \epsilon}, \quad j \in \{1, 2, 3, 4\}
\end{equation}

\textbf{4) Radial MCP Geometry ($\mathbf{f}_{15\dots19}$):} Distances from wrist to metacarpophalangeal (MCP) joints normalized by $D_{\text{palm}}$.

\subsection{Lightweight Neural Classification and Decision Logic}
The 19-dimensional feature vector is processed by a compact Multilayer Perceptron (MLP) consisting of an input layer (19), two hidden layers (64 and 32 neurons with ReLU activation and 20\% dropout), and a 6-node softmax output layer corresponding to the operational vocabulary: \texttt{STOP}, \texttt{GO}, \texttt{LEFT}, \texttt{RIGHT}, \texttt{BACK}, and \texttt{FOLLOW}. 

The centralized \texttt{brain\_node} manages state transitions using a sliding temporal majority-vote buffer of size $W=5$ frames. A gesture command is confirmed only when a single class achieves an agreement threshold $\tau \ge 4$, eliminating single-frame classification flicker. Active subject locking is enforced via a normalized central spatial zone $[x_{\min}, x_{\max}] = [0.20, 0.80]$ in normalized image coordinates, rejecting gestures from peripheral bystanders.

\begin{table}[t]
\centering
\caption{Physical Trial Gesture Recognition Results (180 Total Trials)}
\label{tab:physical_trials}
\begin{tabular}{lcccc}
\toprule
\textbf{Gesture Command} & \textbf{Trials} & \textbf{Success} & \textbf{Accuracy (\%)} & \textbf{Robot Response} \\
\midrule
STOP   & 30 & 27 & 90.00\% & Complete halt ($v_x=0.0$) \\
FOLLOW & 30 & 29 & 96.67\% & Visual servoing buffer \\
GO     & 30 & 28 & 93.33\% & Forward ($v_x=+0.20$\,m/s) \\
LEFT   & 30 & 30 & 100.00\% & Counter-clockwise turn \\
RIGHT  & 30 & 30 & 100.00\% & Clockwise turn \\
BACK   & 30 & 30 & 100.00\% & Reverse ($v_x=-0.15$\,m/s) \\
\midrule
\textbf{Aggregate Total} & \textbf{180} & \textbf{174} & \textbf{96.67\%} & \textbf{Deterministic Control} \\
\bottomrule
\end{tabular}
\end{table}

\section{Experimental Results and Analysis}

\subsection{Empirical Gesture Classification Performance}
Physical testing was conducted with three human operators under daytime natural illumination and nighttime artificial fluorescent lighting across interaction distances of 1.00\,m, 1.75\,m, and 2.50\,m. As shown in Table~\ref{tab:physical_trials}, the system achieved an aggregate recognition accuracy of 96.67\% (174 out of 180 trials). Directional commands achieved 100\% reliability due to pronounced angular feature separation. The three unclassified STOP gestures resulted from acute hand tilt angles exceeding $45^\circ$, which partially occluded palm landmarks.

Across operational distances, accuracy recorded 100.0\% at 1.00\,m (54/54), 97.2\% at 1.75\,m (70/72), and 92.6\% at 2.50\,m (50/54). Illumination invariance was verified with identical 96.7\% accuracy under both daylight and artificial lighting regimes.

\begin{table}[t]
\centering
\caption{End-to-End Latency Budget Breakdown on Embedded Edge Hardware}
\label{tab:latency_budget}
\begin{tabular}{llr}
\toprule
\textbf{Stage} & \textbf{Pipeline Operation} & \textbf{Latency (ms)} \\
\midrule
1 & Monocular V4L2 capture ($640 \times 480$ @ 20 FPS) & 50.0 \\
2 & 21-point hand landmark extraction (MediaPipe) & 38.5 \\
3 & 19D scale-invariant feature extraction & 1.2 \\
4 & MLP neural inference (ONNX / CPU) & 2.8 \\
5 & ROS~2 intra-process and DDS transport & 4.5 \\
6 & ESP32 UART serial bridge and motor acceleration ramp & 35.0 \\
\midrule
\textbf{Total} & \textbf{Sensor-to-Torque Turnaround Latency} & \textbf{132.0} \\
\bottomrule
\end{tabular}
\end{table}

\subsection{End-to-End Latency Budget and Safety Compliance}
Table~\ref{tab:latency_budget} provides the empirical latency budget measured from photon arrival at the camera sensor to motor torque generation. The total round-trip latency measured 132.0\,ms on the Raspberry Pi 5 CPU. This performance satisfies the 150\,ms human-interaction reaction threshold and complies with ISO 15066:2016 safety requirements for collaborative mobile robots. Vehicle braking tests confirmed complete physical halting within 98\,ms following command issuance, guaranteeing safe stopping distances during close-proximity human interactions.

\subsection{Resource Utilization and Multi-Node Throughput}
Under concurrent execution of YOLOv8n, MediaPipe HandLandmarker, the MLP classifier, an 8-frame LSTM trajectory predictor, and SLAM Toolbox, the quad-core Cortex-A76 CPU stabilized at 311\% utilization ($\approx 78\%$ per core). Active physical memory consumption peaked at 4.2\,GB, preserving over 3.8\,GB of unallocated LPDDR4X headroom on the 8\,GB hardware platform, preventing out-of-memory kernel terminations. Measured node rates sustained 20.0\,Hz for camera streaming, 10.0\,Hz for gesture publishing, and 12.6\,Hz for planar LiDAR mapping.

\subsection{Comparative Benchmarking}
Table~\ref{tab:benchmarking} benchmarks the proposed system against state-of-the-art literature. Unlike prior frameworks requiring specialized external depth hardware or validated strictly in simulation, our system achieves competitive physical accuracy (96.67\%) and low latency (132.0\,ms) natively on an embedded CPU with closed-loop SLAM mobile navigation.

\begin{table}[t]
\centering
\caption{Comparative Benchmarking Against Prior Intention Frameworks}
\label{tab:benchmarking}
\resizebox{\columnwidth}{!}{
\begin{tabular}{lcccc}
\toprule
\textbf{Framework} & \textbf{Sensor} & \textbf{Compute} & \textbf{Accuracy} & \textbf{Mobility / SLAM} \\
\midrule
Tsitos et al. \cite{tsitos2022} & RGB-D & Desktop GPU & 91.2\% & Stationary UR3 \\
Mahmud et al. \cite{mahmud2022} & Kinect 3D & Workstation & 94.4\% & Simulation Only \\
Li et al. \cite{li2023} & 2D LiDAR+Cam & High-End PC & N/A & Isaac Sim Only \\
Muhtadin et al. \cite{muhtadin2025} & RGB Monoc & Jetson Nano & 95.8\% & Stationary UR5 \\
\textbf{Proposed System} & \textbf{RGB Monoc} & \textbf{Raspberry Pi 5} & \textbf{96.67\%} & \textbf{Physical Mobile SLAM} \\
\bottomrule
\end{tabular}
}
\end{table}

\section{Conclusion}
This paper presented an edge-feasible human intention recognition and autonomous navigation system deployed on a mobile robot without GPU acceleration. By mathematically decoupling hand features from distance and scale, the system sustained 96.67\% physical recognition accuracy across varied operational distances and illumination conditions. With an empirical latency of 132.0\,ms, the system provides safe, responsive human-robot interaction compliant with ISO 15066. Future work will extend this framework to multi-robot collaborative coordination and embedded tensor-accelerated NPU deployment.

\section*{Acknowledgment}
This work was conducted at the Department of Computer Engineering, Faculty of Engineering, Ghana Communication Technology University (GCTU), Accra, Ghana.

\bibliographystyle{IEEEtran}
\bibliography{references}

\end{document}
